\chapter{Introduction}
\label{chap:intro}

This chapter states the problem, identifies the gap in existing work that
motivates the project, sets out the aims and objectives against which the work
is assessed, and describes the structure of the remainder of the report.

\section{Background}
\label{sec:background}

Network intrusion detection systems classify traffic as benign or as belonging
to a class of attack. How well they work depends on the breadth of attack
traffic they have been trained on, which creates a difficulty: an organisation
would benefit from learning from attacks other organisations have seen, but
network traffic is among the most sensitive data it holds. Flow records reveal
internal addressing, service topology, user behaviour, and the fact of a
compromise when one has occurred.

Federated learning \parencite{mcmahan-2017} offers a way around this. Each
participant trains on its own records and transmits only model updates, which a
server averages into a shared model. No raw traffic leaves the participant.

The difficulty is that averaging works well when participants hold similar
data and poorly when they do not. An organisation running web services sees
different attacks from one running industrial control systems. Under this kind
of statistical heterogeneity a single averaged model fits nobody well, an
effect documented across the federated learning literature and particularly
sharp in intrusion detection, where attack classes are rare, unevenly
distributed, and specific to the environment that produced them.

Personalised federated learning addresses this by giving each participant a
model of its own while still sharing what can be shared. Multi-tier variants go
further and introduce an intermediate level between the device and the server,
on the argument that participants fall into groups and that a group-level model
captures what its members have in common without forcing agreement across
groups that have little to share.

\section{The base method}
\label{sec:base-method}

This project takes PerMFL \parencite{permfl-2024} as its base. PerMFL is a
three-tier personalised federated learning method with a device model, a team
model, and a global model, coupled by squared Euclidean penalties. It provides
convergence guarantees, linear for smooth strongly convex problems and
sub-linear for smooth non-convex ones, and its authors published an
implementation together with seven baseline algorithms.

Three properties made it the right base. It is reproducible, which most of the
adjacent literature is not. It carries a convergence proof, which the clustered
federated learning methods it competes with do not. And its stated premise, that
it applies ``when there are known team structures across devices'', maps onto
intrusion detection, where attack families and network segments provide exactly
that kind of structure.

\section{Problem statement}
\label{sec:problem}

PerMFL has been evaluated only on image benchmarks and a synthetic tabular
dataset. Intrusion detection data differs from those in ways that bear directly
on the method's assumptions. Classes are severely imbalanced, with benign
traffic accounting for over eighty per cent of the CICIDS2017 capture. Several
attack classes have double-digit sample counts against millions of benign
flows. Traffic arrives in time order and attacks occur in contiguous bursts, so
the split between training and evaluation data is a methodological decision
rather than a formality. None of these conditions is present in MNIST.

The project asks whether the method transfers, what has to change for it to
transfer, and whether the architectural claim that motivates it holds on this
kind of data.

\section{Aims and objectives}
\label{sec:aims}

The aim is to establish whether PerMFL is suitable for network intrusion
detection and, where it is not, to identify why and what corrects it.

\begin{enumerate}
  \item Reproduce the published result on the authors' own dataset, to
        establish that the implementation under test behaves as the paper
        describes.
  \item Build a leakage-free data pipeline for CICIDS2017, NSL-KDD and
        TON-IoT, with a documented partitioning scheme and an evaluation
        protocol appropriate to imbalanced multi-class detection.
  \item Measure the method under controlled variation of client heterogeneity,
        reporting each result against the floor a trivial classifier achieves
        and the ceiling the same model reaches without federation.
  \item Test the architectural premise directly by varying the team structure
        the method is built around.
  \item Where the method underperforms, identify the mechanism and evaluate a
        correction.
\end{enumerate}

\section{Contributions}
\label{sec:contributions}

The work makes four contributions.

An exact reproduction of the published EMNIST-10 result, matching the reported
personalised accuracy to within 0.04 percentage points, together with sixteen
documented defects in the released implementation. One of these, an argument
binding error, silently overrides a tunable hyperparameter with an unrelated
loop count.

Evidence that the team tier, the architectural feature that distinguishes
PerMFL from two-tier methods, contributes almost nothing within the parameter
region its own convergence conditions permit. Removing the tier entirely
changes personalised macro F1 by 0.0002 on CICIDS2017 and by 0.0000 on the
authors' own EMNIST-10 setup.

An explanation of why, derived from the update rule rather than inferred from
results. A single parameter governs both the device-to-team penalty and the
weight the team model places on its members, and the values these two roles
require are opposed.

A correction that separates them, evaluated across three datasets and four
partitioning schemes with five paired random seeds. It improves the global
model on three of four partitions and reduces its variance across seeds by two
to three orders of magnitude in all four, with one regression that is reported
rather than omitted.

\section{Report structure}
\label{sec:structure}

Chapter~\ref{chap:litreview} reviews personalised and clustered federated
learning, federated intrusion detection, and the data quality problems specific
to CICIDS2017. Chapter~\ref{chap:requirements} states the functional and
non-functional requirements and the ethical and legal position.
Chapter~\ref{chap:design} describes the experimental design, the partitioning
schemes and the evaluation protocol. Chapter~\ref{chap:implementation} covers
the data pipeline, the metric instrumentation and the modifications to the base
implementation. Chapter~\ref{chap:testing} sets out how the pipeline was
verified. Chapter~\ref{chap:evaluation} presents the results and evaluates them
critically. Chapter~\ref{chap:conclusion} concludes and identifies further
work.
